\documentclass[11pt]{article}
\usepackage[margin=1in]{geometry}
\usepackage{booktabs}
\usepackage{hyperref}

\title{Finding Clashes in Coding-Agent Memory:\\Formal Checking, LLM Judges, and Where Each Fails}
\author{CS2520 Project}
\date{December 2026}

\begin{document}
\maketitle

% Topic sentences only (draft skeleton, 2026-11-23). Numbers marked TBD are pending.

\begin{abstract}
Coding agents accumulate memory rules that can contradict one another, and we ask whether
translating each rule into a small formal vocabulary and checking it with Z3 finds these
clashes better than asking an LLM directly.
\end{abstract}

\section{Introduction}
Memory folders and \texttt{CLAUDE.md} files grow by accretion, so an agent can be told both to
force-push a branch and never to force-push it, with no warning to the user.

\section{Related Work}
Prior work on natural-language-to-logic translation shows that solver hybrids win when
translation is narrow and reasoning is deep or spans several requirements, and lose when
understanding the text is the hard part.

\section{What Counts as a Clash}
We define a clash as a situation, named by the rules, in which no single action follows all of
them, reading an explicit ``may'' as a strong permission.

\section{Method}
\subsection{Vocabulary}
A fixed vocabulary of git, file, and task actions with conditions lets every rule in scope be
written as a small set of obligation, prohibition, permission, and conditional clauses.
\subsection{Two-Step Translation}
An LLM first lists each rule's requirements with exact quotes and scopes, then encodes them,
and a code check makes sure every requirement is either encoded or explicitly left out.
\subsection{Checking Pairs and Groups with Z3}
Z3 searches every situation the rules distinguish for one where no action satisfies them all,
for pairs and for groups of three rules, and returns that situation as a witness.

\section{Data}
We collected 1,287 rules from five real memory folders, labelled translation sets with two
independent labellers, and built planted and natural clash sets.

\section{Experiments and Results}
\subsection{Translation Accuracy}
Two-step translation beats one-step translation by 14 points on in-scope rules when either
labeller's defensible reading counts, while the two labellers themselves agree on only 67\%.
\subsection{Pairwise Clashes}
On planted git/file clashes the checker finds 15--18 of 20 with no false alarms, but a direct
LLM judge finds all of them, so on single pairs the LLM is the stronger baseline.
\subsection{Clashes Among Three Rules}
On groups of three rules that clash only together, the set-level checker finds only the group
whose rules all translate (1 of 8, no false alarms), while an LLM reading the whole folder finds
5 of 8, so translation and vocabulary, not set-level reasoning, decide the outcome.
\subsection{Real Folders at Full Size}
Across five full folders, most real clashes concern tool choice rather than git or files, and a
single whole-folder LLM pass misses many of them.

\section{Discussion}
Z3 is exact about what it receives, so the failures we observe come from translating rules one
at a time and from choosing which rules to compare, not from the reasoning.

\section{Limitations}
Our labels come from model labellers on five folders, recall over all pairs is unknown, and the
development set was partly tuned on.

\section{Conclusion}
Formal checking is worth keeping for clashes inside its vocabulary, especially those that need
several rules at once, and an LLM judge should handle the rest with a person settling ambiguous
cases.

\end{document}
